% Sections 1-3 of: The Information-Quality Paradox in On-Policy Distillation
% NeurIPS 2026 submission
% Requires: neurips_2026.sty, booktabs, graphicx, amsmath, amssymb

\section{Introduction}
\label{sec:intro}

On-policy knowledge distillation has become the dominant paradigm for compressing reasoning capabilities from large language models into smaller ones~\citep{agarwal2024gkd, gu2024minillm}.
The standard recipe is straightforward: the student generates responses on its own distribution, the teacher scores them via a forward pass, and the student is trained to minimize the reverse KL divergence at every token position.
An implicit assumption underlies this approach---that all token positions contribute equally useful supervision.
We show this assumption is wrong, and violating it is not merely wasteful but actively harmful.

We identify an \emph{information-quality paradox}: the tokens with the highest KL divergence between teacher and student---seemingly the most informative for learning---produce the \emph{worst} distillation outcomes.
Concretely, selecting the 100 tokens per response with the highest KL divergence captures 93\% of the total KL signal yet achieves only 53.4\% on MATH-500, while the first 100 positional tokens capture only 46\% of total KL yet achieve 65.9\% (Table~\ref{tab:token_selection}).
This 12.5 percentage-point gap inverts the information-theoretic intuition that concentrating on the largest gradients should accelerate learning.

The resolution lies in what high-KL tokens actually encode.
Per-token analysis over two million tokens reveals that the highest-KL positions are dominated by format and style disagreements---LaTeX delimiters (\texttt{\textbackslash(}, \texttt{\textbackslash[}) with KL exceeding 10 nats per token, planning phrases (``Therefore'', ``First'') with KL $\approx 1.75$---rather than mathematical reasoning content, where KL is an order of magnitude lower ($\approx 0.28$ for numeric tokens).
Meanwhile, teacher signal quality decays monotonically with position: mean KL drops from 1.9 to 0.4, teacher entropy falls from 0.26 to 0.17, and teacher-student agreement rises from 75\% to 89\% across the response.
Late-position supervision is the teacher rubber-stamping the student's choices, not correcting its reasoning.

Based on this analysis, we propose \emph{positional distillation}: computing loss only on the first $N$ response tokens.
This requires changing a single line of code (\texttt{loss\_mask[:, N:] = 0}), yet delivers a 24$\times$ wall-clock speedup, 4$\times$ memory reduction, and matches or exceeds full-sequence performance.
We validate positional distillation across three model families (Qwen, Gemma, Llama), three tasks (mathematics, function calling, coding), and two training methods (LoRA, full fine-tuning), with multi-seed evaluation confirming stability.
On function calling, the distilled student exceeds the teacher by up to 9.7 percentage points---evidence that the method extracts clean reasoning signal while filtering teacher output noise.

Full-sequence distillation, by contrast, is fragile: it collapses in 2 out of 3 seeds on Qwen math (mean 56.8\% $\pm$ 6.1 vs.\ 62.4\% $\pm$ 2.9 for positional), drops below baseline on Gemma, and degrades from 55.3\% to 32.0\% on Llama function calling.

\begin{figure}[t]
    \centering
    % Three-panel teaser figure
    % (a) Token selection paradox: bar chart, avg@4 vs strategy, annotated with KL coverage %
    % (b) Per-position KL curves for Qwen, Gemma, Llama (different decay profiles)
    % (c) Cross-family: grouped bars of baseline / pos-best / fullseq for each family
    \includegraphics[width=\textwidth]{figures/figure1_teaser.pdf}
    \caption{\textbf{The information-quality paradox.}
    \textbf{(a)}~Token selection comparison on MATH-500 (K=100 tokens). Selecting by highest KL (93\% coverage) performs worst; the contiguous prefix (46\% coverage) performs best.
    \textbf{(b)}~Per-position mean KL across three model families. Qwen and Gemma are front-loaded; Llama is flat. Dashed lines mark the optimal position limit for each.
    \textbf{(c)}~Cross-family math results. Positional distillation consistently outperforms full-sequence, which collapses or degrades in every setting.}
    \label{fig:teaser}
\end{figure}

Our contributions are:
\begin{enumerate}[leftmargin=*,itemsep=2pt,topsep=2pt]
    \item \textbf{Information-quality paradox.} We demonstrate that high-KL tokens in on-policy distillation correspond to format and style noise, not reasoning quality. Training on them produces the worst outcomes (\S\ref{sec:paradox_result}), overturning the intuition that the largest divergences are the most informative.
    \item \textbf{Signal quality analysis.} Per-position analysis across three model families reveals that teacher signal quality decays with position. Cross-model KL profiles predict optimal position limits (\S\ref{sec:cross_model}).
    \item \textbf{Cross-family, cross-task validation with multi-seed.} Positional distillation matches or exceeds full-sequence across 3 model families $\times$ 3 tasks $\times$ 2 training methods, with 24$\times$ speedup and 4$\times$ memory reduction. Multi-seed runs confirm 2.6$\times$ lower variance than full-sequence.
    \item \textbf{Cascade effect.} Early-token distillation changes entire generation trajectories: KL reduction extends 35--38\% beyond the trained range, and late tokens change more than early tokens.
\end{enumerate}

\paragraph{Relationship to concurrent work.}
\citet{zhang2025fast} independently discover that prefix truncation accelerates on-policy distillation.
Our work differs in three respects: we explain \emph{why} prefix tokens are superior (signal quality decay, the information-quality paradox), demonstrate that positional distillation is \emph{better} and not merely faster (via the token selection comparison), and validate across three model families and three tasks.
\citet{li2025rethinking} show that teacher-student vocabulary overlap decreases with depth, complementing our findings from a different analytical angle.


%==============================================================================
\section{Background and Related Work}
\label{sec:background}

\subsection{On-Policy Knowledge Distillation}

Knowledge distillation transfers capabilities from a teacher model to a smaller student by training on soft targets~\citep{hinton2015distilling}.
For autoregressive language models, sequence-level KD~\citep{kim2016sequence} generates from the teacher and trains the student on these sequences.
On-policy variants reverse this: the student generates, the teacher scores, and the student trains on the divergence between distributions~\citep{agarwal2024gkd, gu2024minillm}.

Let $p_\theta$ denote the student and $q_\phi$ the teacher.
Given a prompt $x$, the student generates a response $y_{1:T} \sim p_\theta(\cdot \mid x)$.
The standard on-policy reverse KL objective is:
\begin{equation}
\mathcal{L}_{\mathrm{full}} = \mathbb{E}_{y \sim p_\theta}\left[\sum_{t=1}^{T} D_{\mathrm{KL}}\!\big(q_\phi(\cdot \mid x, y_{<t}) \;\|\; p_\theta(\cdot \mid x, y_{<t})\big)\right].
\label{eq:full_kl}
\end{equation}
This objective treats all $T$ response positions uniformly.
Reverse KL is preferred over forward KL in the on-policy setting because it avoids the mode-averaging behavior that arises when the student cannot cover the full teacher distribution~\citep{gu2024minillm}.

\subsection{Token Selection in Training}

Several lines of work question whether all tokens deserve equal loss weight.
In reinforcement learning for LLMs, \citet{wang2025beyond} show that entropy-based token selection improves RLHF by focusing gradients on high-uncertainty positions.
Rho-1~\citep{lin2024rho1} selects tokens by reference-model loss for pretraining.
SelecTKD~\citep{chen2025selectkd} and AdaKD~\citep{wu2025adakd} apply adaptive weighting in offline distillation.

Most relevant to our work are two concurrent papers.
\citet{zhang2025fast} propose prefix truncation for on-policy distillation, framing it as an efficiency heuristic that preserves 90\% of performance at 3$\times$ speedup.
They do not investigate \emph{why} prefix tokens are superior, nor compare against information-theoretic selection strategies.
\citet{li2025rethinking} analyze teacher-student vocabulary overlap and show that the teacher's reward signal degrades with response depth.
Their analysis complements ours but they do not propose position-restricted loss or validate across model families.

\subsection{Our Position}

We bridge these threads by providing the scientific explanation for why prefix truncation works: teacher signal quality decays with position, and high-KL tokens are noise rather than signal.
Our token selection experiment (Table~\ref{tab:token_selection}) directly compares positional, entropy-based, KL-based, and random strategies under matched compute, revealing the information-quality paradox.
Cross-family validation across Qwen, Gemma, and Llama demonstrates that the finding is not architecture-specific.


%==============================================================================
\section{The Information-Quality Paradox}
\label{sec:paradox}

This section presents the core analytical finding: in on-policy distillation, selecting tokens by information content (KL divergence) produces systematically worse outcomes than a simple positional prefix.
We first establish the paradox empirically, then resolve it by examining what high-KL tokens encode.

\subsection{Experimental Setup}
\label{sec:paradox_setup}

We study the Qwen family: student Qwen2.5-Math-1.5B, teacher Qwen3-1.7B, on MATH-500~\citep{hendrycks2021math}.
The student generates full-length on-policy responses (mean length 386 tokens), and we select $K{=}100$ tokens per response for loss computation under different strategies.
All strategies generate full-length sequences and differ only in which tokens receive gradient signal, ensuring identical compute.
Training uses LoRA ($r{=}32$, $\alpha{=}64$) with learning rate $5{\times}10^{-5}$, processing 3{,}200 problems over 200 steps.
Evaluation reports avg@4 on MATH-500 with temperature 0.7.

The selection strategies are:
\textbf{prefix}---first 100 response tokens;
\textbf{top-kl}---100 tokens with highest per-token KL divergence;
\textbf{ent-teacher}/\textbf{ent-student}---100 tokens with highest teacher/student entropy;
\textbf{random}---100 uniformly sampled tokens;
\textbf{middle}/\textbf{last}---100 contiguous tokens centered at the midpoint or end of the response.
For each strategy, we compute its \emph{KL coverage}: the fraction of total per-response KL captured by the selected tokens.

\subsection{The Paradox: High-KL Tokens Perform Worst}
\label{sec:paradox_result}

\begin{table}[t]
\centering
\caption{\textbf{Token selection comparison} (MATH-500, avg@4, $K{=}100$ tokens selected per response). KL coverage measures the fraction of total response KL captured. Higher coverage does not produce better outcomes---the relationship is inverted.}
\label{tab:token_selection}
\small
\begin{tabular}{@{}lccl@{}}
\toprule
Strategy & KL Coverage & Best avg@4 & Stability \\
\midrule
\textbf{prefix-100}  & \textbf{45.6\%} & \textbf{65.85\%} & Stable across all steps \\
ent-teacher       & 50.5\% & 63.30\% & Mild degradation \\
ent-student       & 67.0\% & 62.70\% & Stable \\
middle-100        & 15.8\% & 61.80\% & Stable \\
random-100        & 21.1\% & 61.25\% & Collapses at step 150 \\
last-100          & 13.9\% & 62.30\% & Collapses at step 150 \\
\textbf{top-kl-100} & \textbf{93.2\%} & \textbf{53.35\%} & Below baseline throughout \\
\bottomrule
\end{tabular}
\end{table}

Table~\ref{tab:token_selection} presents the central result.
The top-KL strategy captures 93.2\% of total KL divergence---nearly all of the information-theoretic signal---yet achieves only 53.35\%, which is 2.4 percentage points \emph{above} the no-distillation baseline of 50.95\% and 12.5 points below the prefix strategy.
The prefix, despite capturing less than half the KL signal (45.6\%), achieves the highest accuracy (65.85\%) and remains stable across all training steps.

The correlation between KL coverage and performance is \emph{negative}: strategies that capture more KL divergence perform worse.
This inverts the standard intuition from curriculum learning and active learning that concentrating on the largest errors accelerates learning.

Two additional patterns emerge.
First, \emph{contiguity matters}: the middle-100 strategy (15.8\% KL coverage) outperforms random-100 (21.1\% coverage) and avoids its collapse, suggesting that training on contiguous token spans provides more coherent gradients than scattered positions.
Second, \emph{position matters directionally}: strategies that include late-position tokens (random, last) exhibit catastrophic collapse at step 150, while prefix and middle remain stable.

\subsection{Resolution: What High-KL Tokens Actually Encode}
\label{sec:resolution}

To understand why high-KL tokens harm distillation, we classify over two million tokens from on-policy trajectories into semantic categories and measure mean KL by category.

\begin{table}[t]
\centering
\caption{\textbf{Mean KL divergence by token category.} High-KL tokens are dominated by format and style disagreements, not mathematical reasoning. Token counts are from 59{,}936 on-policy trajectories.}
\label{tab:token_classification}
\small
\begin{tabular}{@{}llc@{}}
\toprule
Category & Example tokens & Mean KL \\
\midrule
math\_latex & \texttt{\textbackslash(}, \texttt{\textbackslash[}, \texttt{\textbackslash\textbackslash} & 3.99 \\
planning & ``To'', ``Therefore'', ``First'' & 1.75 \\
structural & \texttt{**}, \texttt{:} & 0.89 \\
math\_operator & $=$, $+$, $-$ & 0.38 \\
math\_number & 0--9 & 0.28 \\
\bottomrule
\end{tabular}
\end{table}

Table~\ref{tab:token_classification} reveals a clear pattern: the tokens carrying the highest KL are \emph{format tokens}---LaTeX delimiters, structural markers, and discourse connectives---not tokens encoding mathematical content.
Individual LaTeX tokens reach KL values of 10--13 nats; stylistic tokens like ``Please'' and ``Sure'' exceed 25 nats.
These reflect disagreements about \emph{how} to present answers, not \emph{what} the answer is.

The mechanistic explanation follows from the structure of KL divergence.
At positions where the teacher is uncertain (high teacher entropy) but disagrees with the student's choice, KL is large.
Direct measurement confirms this: top-KL tokens have teacher entropy of 0.449, nearly double the 0.235 of prefix tokens~(Table~\ref{tab:entropy_decomposition}).
The ideal distillation signal---low teacher entropy (confident correction) with high KL (student is wrong)---is concentrated at early positions.
The top-KL strategy instead selects for \emph{confused disagreement}: both models are uncertain, and the gradient signal is noise.

\begin{table}[t]
\centering
\caption{\textbf{Teacher entropy decomposition.} Top-KL tokens correspond to positions where the teacher is confused (high entropy), not where it has superior knowledge.}
\label{tab:entropy_decomposition}
\small
\begin{tabular}{@{}lccc@{}}
\toprule
Token set & Teacher entropy & Student entropy & Mean KL \\
\midrule
Top-KL ($K{=}100$) & 0.449 & 0.696 & 4.09 \\
Prefix (0--99) & 0.235 & 0.466 & 2.02 \\
Late (100+) & 0.158 & --- & 1.81 \\
\bottomrule
\end{tabular}
\end{table}

\subsection{Per-Position Signal Quality Decay}
\label{sec:signal_decay}

The token classification explains the paradox at the individual-token level.
We now examine the positional structure that gives rise to it.
Table~\ref{tab:signal_decay} reports three complementary indicators of teacher signal quality as a function of response position, measured over 59{,}936 on-policy trajectories.

\begin{table}[t]
\centering
\caption{\textbf{Per-position signal quality} (Qwen2.5-Math-1.5B $\to$ Qwen3-1.7B). All three indicators show monotonic decay: the teacher becomes less informative, less certain, and more confirmatory with position.}
\label{tab:signal_decay}
\small
\begin{tabular}{@{}lcccc@{}}
\toprule
Position range & Mean KL & Teacher entropy & Agreement rate & $N$ tokens \\
\midrule
0--50      & 1.91 & 0.26 & 75.4\% & 2{,}467 \\
50--100    & 0.94 & 0.21 & 81.9\% & 2{,}450 \\
100--200   & 0.65 & 0.16 & 86.3\% & 4{,}844 \\
200--300   & 0.52 & 0.16 & 87.8\% & 4{,}497 \\
300--500   & 0.43 & 0.17 & 88.8\% & 8{,}124 \\
\bottomrule
\end{tabular}
\end{table}

Three indicators paint a consistent picture:
\begin{itemize}[leftmargin=*,itemsep=2pt,topsep=2pt]
    \item \textbf{KL divergence} drops by 4.4$\times$ from the first 50 tokens (1.91) to positions 300--500 (0.43). The teacher's ``disagreement budget'' is front-loaded.
    \item \textbf{Teacher entropy} declines from 0.26 to 0.17. Late-position teacher distributions are sharply peaked, offering little distributional information beyond the top-1 prediction.
    \item \textbf{Agreement rate} rises from 75.4\% to 88.8\%. At late positions, the teacher's argmax matches the student's generated token nearly nine times out of ten---supervision is confirmatory rather than corrective.
\end{itemize}

This decay has a natural interpretation in mathematical reasoning.
Early tokens encode the problem-solving strategy: which technique to apply, how to decompose the problem, what intermediate quantities to compute.
These are precisely the positions where the teacher's superior reasoning capability manifests as a different---and often better---distributional preference.
Late tokens encode arithmetic execution and answer formatting, where teacher and student largely agree and the remaining disagreement is stylistic.

\subsection{Cross-Model KL Profiles Predict Optimal Position}
\label{sec:cross_model}

If signal quality decay drives the paradox, different model families with different KL profiles should require different position limits.
We measure pre-training KL profiles for two additional families: Gemma (gemma-2-2b-it $\to$ gemma-3-4b-it) and Llama (Llama-3.2-1B $\to$ Llama-3.1-8B).

\begin{table}[t]
\centering
\caption{\textbf{Cross-model KL profiles and optimal position limits.} Despite very different profile shapes, the optimal position limit consistently falls at $\sim$40--50\% cumulative KL coverage. Flat-KL families require more tokens.}
\label{tab:cross_model_kl}
\small
\begin{tabular}{@{}lcccc@{}}
\toprule
Model pair & KL ratio (first 100 / rest) & Profile shape & Optimal $N$ & Coverage at optimum \\
\midrule
Qwen 1.5B $\to$ 1.7B & 2.81$\times$ & Front-loaded & 100 & ${\sim}$44\% \\
Llama 1B $\to$ 8B    & 1.15$\times$ & Flat         & 200 & ${\sim}$45\% \\
\bottomrule
\end{tabular}
\end{table}

Table~\ref{tab:cross_model_kl} reveals a striking regularity.
The Qwen pair has a sharply front-loaded KL profile (2.81$\times$ more KL in the first 100 tokens than the rest), and its optimal position limit is $N{=}100$.
The Llama pair has a nearly flat profile (1.15$\times$ ratio), and its optimal position limit is $N{=}200$---roughly twice as many tokens are needed to capture the same fraction of useful signal.
Despite these very different shapes, the optimal position in both cases falls at approximately 40--50\% of cumulative KL coverage.

\begin{figure}[t]
    \centering
    % Two-panel analysis figure
    % (a) Per-position KL curves for Qwen and Llama, with optimal N marked
    % (b) Cumulative KL fraction vs position, with 40-50% coverage band shaded
    \includegraphics[width=\textwidth]{figures/figure2_kl_profiles.pdf}
    \caption{\textbf{KL profiles and the 40--50\% coverage rule.}
    \textbf{(a)}~Per-position mean KL for Qwen (front-loaded) and Llama (flat). Dashed lines indicate the empirically optimal position limit.
    \textbf{(b)}~Cumulative KL coverage as a function of position limit. Despite different profile shapes, the optimal $N$ for both families falls in the 40--50\% cumulative KL band (shaded).}
    \label{fig:kl_profiles}
\end{figure}

This suggests a practical heuristic: to select $N$ for a new model pair without a grid search, generate a small set of on-policy trajectories (${\sim}$1{,}000), compute the per-position KL profile, and choose $N$ at the 40--50\% cumulative KL threshold.
Front-loaded profiles (Qwen-like) will yield $N \approx 50$--$100$; flat profiles (Llama-like) will yield $N \approx 150$--$200$.
We validate this heuristic in the experiments (\S\ref{sec:experiments}).
